\section{Introduction}

\subsection{Problem Statement and Motivation}

Epilepsy, a chronic neurological disorder affecting 60 million people worldwide,
with approximately one-third of patients remaining drug-resistant \parencite{husseinFocalNonFocalEpilepsy2018}.
\textcite{sveinssonClinicalRiskFactors2020} found that 69\% of deaths due to a generalized
tonic-clonic seizure could have been prevented if not left unattended, highlighting the demand
for reliable and timely prediction.

While reliable prediction is already possible through electroencephalography (EEG),
its operational complexity renders it unsuitable for ambulatory use.
To address the need for ambulatory detection many are experimenting with
utilizing autonomic and motion signals to detect and predict seizures.

Translating and analysing these often noisy and complex time-series physiological
signals has pushed traditional machine learning (ML) approaches to their limits.
These approaches require manual engineering such as cleaning and extracting of the data,
often failing to capture the complex, non-linear temporal dynamics of these
signals. It is this limitation that deep learning (DL) models are designed to overcome, with their ability of hierarchical and learned feature
extraction uniquely addressing these limitations.

This review therefore examines and assesses recent advances in time-series
deep learning - particularly convolutional neural networks (CNNs)
and long short term memory (LSTMs) -
applied to seizure detection and prediction using non-EEG wearable biosignals.
Central to this analysis is the question of how architectural choices, training paradigms
(e.g., transfer learning, personalization), and thorough prospective evaluation
can converge to meet the stringent clinical requirements of high sensitivity
and minimal false alarm rates (FAR) in real-world deployment.

\subsection{Key Research Question}
How do different deep learning architectures and biosignal modalities excluding EEG
compare in their ability to achieve an optimal trade-off between 
sensitivity and false alarm rate for ambulatory seizure monitoring?